% docs/manual/developer/chapters/11-contributing.tex
% 第 11 章：贡献流程

\chapter{贡献流程}

本章总结向 \openbench{} 主仓库贡献代码的完整流程：从准备改动到 PR 合入到 release。

\section{贡献前：选要做什么}

三类常见贡献：

\begin{itemize}
  \item \textbf{修 bug}：从 GitHub Issues 选 \texttt{good first issue} 或自己遇到的
  \item \textbf{加新功能}：先开 Issue 讨论，避免无用功
  \item \textbf{改文档}：低门槛入门，欢迎从 typo 修复开始
\end{itemize}

\section{Fork → Branch → PR}

标准流程：

\begin{minted}{bash}
# 1. fork on GitHub, then clone
git clone https://github.com/<YOU>/OpenBench.git
cd OpenBench
git remote add upstream https://github.com/CoLM-SYSU/OpenBench.git

# 2. 同步 upstream
git fetch upstream && git checkout main && git merge upstream/main

# 3. 开 feature branch
git checkout -b fix/issue-123-some-bug

# 4. 写代码 + 测试

# 5. ruff 干净 + 测试通过
ruff check src tests && ruff format --check src tests && pytest

# 6. commit
git add <changed files>
git commit -m "fix: short description (#123)"

# 7. push
git push origin fix/issue-123-some-bug

# 8. 在 GitHub 上开 PR
\end{minted}

\section{Commit 信息约定}

参考 \href{https://www.conventionalcommits.org}{Conventional Commits}：

\begin{itemize}
  \item \texttt{feat:} 新功能
  \item \texttt{fix:} 修 bug
  \item \texttt{docs:} 文档
  \item \texttt{refactor:} 不改行为的重构
  \item \texttt{test:} 测试
  \item \texttt{chore:} 杂项（依赖更新、CI 调整）
  \item \texttt{build:} 构建相关
\end{itemize}

\subsection{Scope}

加 scope 表明影响范围：

\begin{itemize}
  \item \texttt{feat(cli):}
  \item \texttt{fix(config):}
  \item \texttt{docs(user):} / \texttt{docs(dev):}
\end{itemize}

\subsection{消息格式}

\begin{minted}{text}
<type>(<scope>): <subject>

<body — optional, multi-line>

<footer — optional, e.g., "Closes #123">
\end{minted}

示例（取自仓库）：

\begin{itemize}
  \item \texttt{fix(cli): init writes flat reference mapping (loader compat)}
  \item \texttt{docs(user): write chapter 1 — overview and installation}
  \item \texttt{feat(manual): add latex\_escape utility with TDD}
\end{itemize}

\section{PR 流程}

\subsection{PR 模板}

\openbench{} 的 PR 描述应包含：

\begin{itemize}
  \item \textbf{What}：改了什么（一句话总结）
  \item \textbf{Why}：为什么（链接到 Issue 或描述场景）
  \item \textbf{How}：怎么改（重点设计决定，不要复述 diff）
  \item \textbf{Testing}：怎么验证（手动 / 自动测试）
  \item \textbf{Closes}：\texttt{Closes \#123}
\end{itemize}

\subsection{Review checklist}

reviewer 会看：

\begin{itemize}
  \item Tests pass + 新代码有测试
  \item ruff lint + format 干净
  \item commit 消息符合约定
  \item 文档更新（CHANGELOG + 相关章节）
  \item 没引入无关改动
  \item 设计与 spec 一致（如果有 spec）
\end{itemize}

\subsection{Review 来回}

\begin{itemize}
  \item Reviewer 留 comment → contributor 修 → push 到同 branch
  \item 不要 force-push（除非 reviewer 要求）
  \item Approved 后由 maintainer 合并
\end{itemize}

\section{Superpowers 工作流：何时使用}

中大型改动用 \modname{superpowers}：

\begin{enumerate}
  \item Spec：\file{docs/superpowers/specs/YYYY-MM-DD-<topic>-design.md}
  \item Plan：\file{docs/superpowers/plans/YYYY-MM-DD-<topic>.md}
  \item Execute：按 plan 提交 commits
  \item Review：\file{docs/superpowers/reviews/YYYY-MM-DD-<topic>.md}（Bug 与决议）
\end{enumerate}

\subsection{用与不用}

\textbf{用}：

\begin{itemize}
  \item 新功能（新 metric、新 GUI page、新统计模块）
  \item 架构改动（schema 重命名、API 重组）
  \item 大文档项目（如本手册）
\end{itemize}

\textbf{不用}：

\begin{itemize}
  \item 单 commit bug 修复
  \item Typo / docstring 改动
  \item 已有功能加测试
\end{itemize}

\subsection{文档命名}

\begin{itemize}
  \item Spec / Plan / Review 同一 topic 用相同日期 + 简称
  \item 一个项目可能多个 plan（如 manual：infrastructure / generators / volume-user / volume-developer / volume-operations）
\end{itemize}

\section{文档更新}

每次代码改动如影响：

\begin{itemize}
  \item 用户可见行为 → 改用户卷对应章节 + 用户卷附录
  \item 内部架构 → 改开发卷对应章节 + 开发卷附录
  \item HPC / 远程 → 改运维卷
\end{itemize}

参考字段、CLI、metrics 等 \emph{自动生成} 的部分，运行 \cli{cd docs/manual \&\& make generated} 让它们自动刷新。

\section{CHANGELOG.md}

每次 release 前更新 \file{CHANGELOG.md}：

\begin{minted}{markdown}
# Changelog

## [3.0.0a1] - 2026-04-30

### Added
- Unified YAML config format (replaces v2.0 multi-file)
- ...

### Fixed
- init_cmd reference shape mismatch with loader (#123)
- ...

### Changed
- ...
\end{minted}

\section{Release 流程（仅 maintainer）}

\begin{enumerate}
  \item 在 main 上确认所有目标 PR 已合
  \item Bump version：\file{src/openbench/\_\_init\_\_.py:\_\_version\_\_}
  \item 更新 CHANGELOG.md
  \item Tag：\texttt{git tag v3.0.0 \&\& git push --tags}
  \item GitHub Release：从 tag 创建，附 changelog
  \item PyPI：\texttt{python -m build \&\& twine upload dist/*}（需 maintainer 凭据）
\end{enumerate}

\section{社区行为}

\begin{itemize}
  \item Issues 与 PR 中保持友好建设性
  \item 给批评 reviewer 时间响应（通常 2 周）
  \item 不要把 GitHub Issue 当成 chat（用 GitHub Discussions / 邮件）
  \item 中英文都欢迎；维护者用英文为主以便国际协作
\end{itemize}

\section{常用资源}

\begin{itemize}
  \item GitHub: \href{https://github.com/CoLM-SYSU/OpenBench}{CoLM-SYSU/OpenBench}
  \item Documentation: \href{https://openbench.readthedocs.io}{openbench.readthedocs.io}（待发布）
  \item 论文（撰写中）：链接 TBD
\end{itemize}

\section{完结}

到此你已读完开发卷全部 11 章，覆盖架构、开发环境、5 个子系统、运行、GUI、测试、贡献流程。下一步：

\begin{itemize}
  \item 看附录 A-E 做参考查阅
  \item 找一个 \texttt{good first issue} 实战
  \item 对架构有想法 → 写 spec 提 Issue 讨论
\end{itemize}

谢谢你愿意贡献 \openbench{}。
